% !TeX program = xelatex
\documentclass[12pt,a4paper,UTF8,fontset=fandol]{ctexart}
\usepackage{BJUTReport}

% ===== 封面信息：在这里统一填写 =====
\renewcommand{\reporttitle}{通信感知一体化研究报告}
\renewcommand{\reportcollege}{信息科学技术学院}
\renewcommand{\reportmajor}{电子信息工程}
\renewcommand{\reportclass}{260261}
% 每位成员一行；个人报告保留一行，小组报告按需增删。
\renewcommand{\reportmembers}{%
  \member{姚思源}{26026115}
  \member{请填写姓名}{请填写学号}
  \member{请填写姓名}{请填写学号}
}
\renewcommand{\reportteacher}{贾克斌}
\renewcommand{\reportdate}{\today}

\begin{document}
\cover

\begin{abstract}

本文介绍通信感知一体化的应用场景和基本原理，围绕无人机配送、道路交通、室内照明、仓储运输和大型活动五类场景，具体说明通感系统部署的位置特点、使用哪些设备、获取什么信息以及怎样发挥作用。在此基础上，介绍系统的基本结构和技术特点，阐明了通感一体化技术在未来智能无线网络中的发展潜力及其面临的关键问题。

\keywords{通感一体化；应用场景；系统结构；智慧交通；低空感知}
\end{abstract}
\thispagestyle{plain}
\clearpage

\tableofcontents
\thispagestyle{plain}
\clearpage
\pagenumbering{arabic}

\section{引言}

通感一体化是一种让无线设备兼顾通信和环境感知的技术。它既可以收发消息，也能够通过通信所使用的无线信号判断通信对象的位置速度等物理信息。在其中通信负责传送信息，感知负责测量环境，两类能力通过设备、信号和处理平台的协同形成服务。

现有公开试验已经展示了车辆与无人机轨迹感知、低空电子围栏和人群密度分析等用途。通感一体化为无线网络增加了观察环境的能力，但实际效果仍取决于覆盖、遮挡、设备和处理条件。其应用价值应通过具体任务判断，例如能否及时发现偏航无人机、给出有用的路口提醒，或减少会议室灯光空开。

\section{技术概述}

\subsection{技术起源}

通感一体化的思路来自通信与雷达两条技术路线的交汇。通信系统把消息加载到无线信号上，再由接收端还原；雷达系统发射信号并分析目标回波，获得距离、速度和方向等信息。两者都需要天线、无线收发电路和信号处理，具有极高的相似性和部分共用的基础。相关综述将早期探索追溯到 20 世纪 60 年代，例如让雷达脉冲在探测之外传递信息。今天面向基站和无线局域网的通感研究，是这一融合思路在新设备和新需求下的发展。[6]

推动它发展的一个直接原因，是许多应用同时需要连接和环境信息。以无人机配送为例，网络要传送任务和飞行状态，飞行控制平台还要实时观察航线附近的目标和障碍物。仅把终端连接起来，无法自动获得所有未联网对象的位置。因此人们开始考虑在传输数据的基础上增加环境感知能力。[9]

另一个原因是无线频谱资源的利用。随着无线通信业务增长和感知应用普及，两类系统对频谱的需求都在增加。因此人们通过让通信和感知功能共享频谱、硬件平台，甚至共同使用发射的无线信号来减少通信与感知系统对独立频谱资源的需求，提高频谱利用效率。

此外，近年来的设备发展让这种融合更具可行性。较宽的带宽为区分距离提供条件，多天线阵列有助于判断方向，数字处理和计算能力的提升则支持更复杂的观测分析。通信与雷达在硬件和处理方法上的接近，使研究逐步从设备共存，扩展到共享信号和联合设计。[6]

因此，通感一体化的发展可以归结为三个相互联系的推动因素：应用需要同时获得连接与环境信息，工程上希望更有效地使用设备和无线资源，技术进步则提供了实现融合的条件。后续章节将通过具体场景说明这些动因如何转化为系统功能。

\subsection{通信和感知}

判断一个方案时是通信还是感知时，可以先问：环境信息从哪里来。通信传递的是发送者准备好的信息，例如文字、视频、控制指令和交通消息。感知则是通过利用无线传播本身来测量环境，是通过分析无线信号的特征来获得物理环境中的测量信息，例如某个方向有没有目标、目标距离多远、是否正在接近。车辆主动报告自己的定位属于通信上传；设备分析车辆反射信号得到距离，则属于无线感知。

两类信息可以互相补充。车辆主动上报能提供身份与状态，但依赖终端连接和数据可信度；信号反射感知可能发现未连接的对象，但未必知道目标身份。实际应用二者相结合可以	提高通信可靠性与传输效率。

\end{document}
